\chapter*{ABSTRACT}
\addcontentsline{toc}{chapter}{ABSTRACT}
\phantomsection

This project demonstrates the design and implementation of an automated hydroponic system allowing soilless growth of crops for efficient resource consumption and improved yield and productivity as compared to conventional farming. Hydroponics refers to the water-based cultivation method in which plants are grown in a solution with calculated nutrient quantities for individual crops. The proposed system consists of sensors such as pH sensor and EC sensor, an ESP32-WROOM32 microcontroller-based control unit, and automated actuators to continuously monitor and regulate key parameters. Automation of the nutrient solution is done through a closed-loop control of pH and EC of the solution where the sensors provide real-time feedback to the controller, which compares measured values with predefined crop-specific standard EC and pH ranges. On the basis of these parameters, dosing pumps are used to add required quantities of acid, base, or concentrated nutrient solutions to compensate for deviations in the nutrient rich solution. Calibration routines are set at periodic time intervals to ensure sensor accuracy and prolonged system reliability. The proposed system is expected to improve pH and EC stability, reduce nutrient wastage and increase crop yield.\\[10pt]
\textit{\textbf{Keywords} --- calibration, closed loop algorithms, dosing, ESP32, EC, hydroponic, pH, TDS}